\documentclass[../../main.tex]{subfiles}
\begin{document}
\section{习题}\label{section:ch04-习题}

\begin{exercise}\label{exercise:ch04-p4-1}
利用包含映射 $\mathbb{H}^n\hookrightarrow\mathbb{R}^n$ 证明: \cref{theorem:ch04-inverse-function} 不能推广到 $M$ 为带边流形的情形。(用于第 80 页。)
\end{exercise}

\begin{exercise}\label{exercise:ch04-p4-2}
设 $M$ 是(不带边的)光滑流形, $N$ 是带边光滑流形, $F:M\to N$ 光滑。证明:若 $p\in M$ 是使 $\mathrm{d}F_p$ 非奇异的点,则 $F(p)\in\operatorname{Int}N$。(用于第 80、87 页。)
\end{exercise}

\begin{exercise}\label{exercise:ch04-p4-3}
为定义域为带边光滑流形的常秩映射,表述并证明秩定理的一个版本。[提示:像在\cref{theorem:ch04-immersion-boundary} 的证明中那样先任意延拓 $F$,然后按秩定理的证明一路推到使用常秩假设的那一步,再说明如何修改延拓后的映射使其具有常秩。]
\end{exercise}

\begin{exercise}\label{exercise:ch04-p4-4}
设 $\gamma:\mathbb{R}\to\mathbb{T}^2$ 为\cref{example:ch04-dense-curve} 中的曲线。证明像集 $\gamma(\mathbb{R})$ 在 $\mathbb{T}^2$ 中稠密。(用于第 502、542 页。)
\end{exercise}

\begin{exercise}\label{exercise:ch04-p4-5}
设 $\mathbb{CP}^n$ 表示 $n$ 维复射影空间,如习题 1-9 所定义。
\begin{enumerate}[(a)]
\item 证明商映射 $\pi:\mathbb{C}^{n+1}\setminus\{0\}\to\mathbb{CP}^n$ 是满射光滑淹没。
\item 证明 $\mathbb{CP}^1$ 微分同胚于 $S^2$。
\end{enumerate}
(用于第 172、560 页。)
\end{exercise}

\begin{exercise}\label{exercise:ch04-p4-6}
设 $M$ 是非空光滑紧流形。证明:对任何 $k>0$,不存在光滑淹没 $F:M\to\mathbb{R}^k$。
\end{exercise}

\begin{exercise}\label{exercise:ch04-p4-7}
设 $M$ 与 $N$ 为光滑流形, $\pi:M\to N$ 是满射光滑淹没。证明:在 $N$ 上不存在任何其他光滑流形结构满足\cref{theorem:ch04-char-submersion} 的结论;换言之,设 $\widetilde{N}$ 表示与 $N$ 相同的集合但可能具有不同的拓扑与光滑结构,并且对每个带边或不带边的光滑流形 $P$,映射 $F:\widetilde{N}\to P$ 光滑当且仅当 $F\circ\pi$ 光滑,证明 $\operatorname{Id}_N$ 是 $N$ 与 $\widetilde{N}$ 之间的微分同胚。[注:这表明\cref{theorem:ch04-char-submersion} 所描述的性质是``特征性的'',其含义与\cref{theorem:appA-quotient-char} 之刻画商拓扑的意义相同。]
\end{exercise}

\begin{exercise}\label{exercise:ch04-p4-8}
本习题说明\cref{theorem:ch04-char-submersion} 的逆命题不成立。设 $\pi:\mathbb{R}^2\to\mathbb{R}$ 由 $\pi(x,y)=xy$ 定义。证明 $\pi$ 满射且光滑,并且对每个光滑流形 $P$,映射 $F:\mathbb{R}\to P$ 光滑当且仅当 $F\circ\pi$ 光滑;但 $\pi$ 不是光滑淹没。
\end{exercise}

\begin{exercise}\label{exercise:ch04-p4-9}
设 $M$ 是连通光滑流形, $\pi:E\to M$ 是拓扑覆叠映射。通过证明 $E$ 上只有一种使 $\pi$ 成为光滑覆叠映射的光滑结构,来完成\cref{proposition:ch04-covering-smooth} 的证明。[提示:利用光滑局部截面的存在性。]
\end{exercise}

\begin{exercise}\label{exercise:ch04-p4-10}
证明\cref{example:ch01-circle-covering}(f) 中定义的映射 $q:S^n\to\mathbb{RP}^n$ 是光滑覆叠映射。(用于第 550 页。)
\end{exercise}

\begin{exercise}\label{exercise:ch04-p4-11}
证明:拓扑覆叠映射是固有的,当且仅当其纤维有限;因此\cref{proposition:ch04-proper-covering} 的逆命题不成立。
\end{exercise}

\begin{exercise}\label{exercise:ch04-p4-12}
利用覆叠映射 $\varepsilon^2:\mathbb{R}^2\to\mathbb{T}^2$(见\cref{example:ch04-smooth-coverings}),证明\cref{example:ch04-submersions-immersions}(d) 中定义的浸入 $X:\mathbb{R}^2\to\mathbb{R}^3$ 可下降为 $\mathbb{T}^2$ 到 $\mathbb{R}^3$ 的光滑嵌入。具体地,证明 $X$ 下降到商上定义了光滑映射 $\widetilde{X}:\mathbb{T}^2\to\mathbb{R}^3$,然后证明 $\widetilde{X}$ 是光滑嵌入,其像即为所给的旋转曲面。
\end{exercise}

\begin{exercise}\label{exercise:ch04-p4-13}
定义映射 $F:S^2\to\mathbb{R}^4$ 为 $F(x,y,z)=\bigl(x^2-y^2,xy,xz,yz\bigr)$。利用\cref{example:ch01-circle-covering}(f) 的光滑覆叠映射与习题 4-10,证明 $F$ 可下降为 $\mathbb{RP}^2$ 到 $\mathbb{R}^4$ 的光滑嵌入。
\end{exercise}

\end{document}